
\documentclass{article}
\usepackage{colortbl}
\usepackage{makecell}
\usepackage{multirow}
\usepackage{supertabular}

\begin{document}

\newcounter{utterance}

\centering \large Interaction Transcript for game `referencegame', experiment `number\_grids', episode 0 with NEW\_Reflect\_Agent.
\vspace{24pt}

{ \footnotesize  \setcounter{utterance}{1}
\setlength{\tabcolsep}{0pt}
\begin{supertabular}{c@{$\;$}|p{.15\linewidth}@{}p{.15\linewidth}p{.15\linewidth}p{.15\linewidth}p{.15\linewidth}p{.15\linewidth}}
    \# & $\;$A & \multicolumn{4}{c}{Game Master} & $\;\:$B\\
    \hline

    \theutterance \stepcounter{utterance}  
    & & \multicolumn{4}{p{0.6\linewidth}}{
        \cellcolor[rgb]{0.9,0.9,0.9}{
            \makecell[{{p{\linewidth}}}]{
                \texttt{\tiny{[P1$\langle$GM]}}
                \texttt{You are given three grids, where each of them is 5 by 5 in size.} \\
\texttt{Grids have empty cells marked with "▢" and filled cells marked with "X".} \\
\texttt{Your task is to generate a referring expression that best describes the target grid while distinguishing it from the two other distractor grids.} \\
\texttt{The first grid is the target grid, and the following two grids are the distractors.} \\
\\ 
\texttt{Target grid:} \\
\\ 
\texttt{□ X X X X} \\
\texttt{□ X □ □ X} \\
\texttt{□ X X X X} \\
\texttt{□ X □ □ X} \\
\texttt{□ X X X X} \\
\\ 
\texttt{Distractor grid 1:} \\
\\ 
\texttt{□ X X X X} \\
\texttt{□ X □ □ X} \\
\texttt{□ X X X X} \\
\texttt{□ □ □ □ □} \\
\texttt{□ X X X X} \\
\\ 
\texttt{Distractor grid 2:} \\
\\ 
\texttt{□ X X X X} \\
\texttt{□ X □ □ X} \\
\texttt{□ X □ □ X} \\
\texttt{□ X □ □ X} \\
\texttt{□ X X X X} \\
\\ 
\texttt{Instruction: Describe the target grid.} \\
\texttt{Generate the referring expression starting with the tag "Expression: " for the given target grid. Omit any other text.} \\
            }
        }
    }
    & \\ \\

    \theutterance \stepcounter{utterance}  
    & \multicolumn{4}{p{0.45\linewidth}}{
        \cellcolor[rgb]{0.95,0.95,0.95}{
            \makecell[{{p{\linewidth}}}]{
                \texttt{\tiny{[P1$|$internal]}}
                \texttt{Reflections from previous episodes:} \\
\texttt{{-}} \\
\texttt{ Failed. My expression "Rows 1, 3, and 5 are fully filled with X's, while rows 2 and 4 have an X only in the first column" did not distinguish the target from Distractor 2, which also has that pattern. I focused on the rows but missed the key detail: in the target, columns 2{-}5 in rows 2 and 4 are empty, but in Distractor 1 those columns have intermittent X's. Next time, I must explicitly note the alternating pattern of *fully* filled rows and rows with only a left{-}edge X, contrasting it with the distractors' internal X's or fully empty columns.} \\
\texttt{{-}} \\
\texttt{ **Reflection:**  } \\
\texttt{I succeeded: my expression distinguished the target from both distractors. I focused on the central X{-}pattern forming a "cross" plus extra Xs on the borders, which distractors lacked. Next time, I would similarly pinpoint unique structural features and avoid vague descriptions.} \\
\texttt{{-}} \\
\texttt{ I succeeded. The target grid has a filled plus sign shape with arms extending to all edges, while both distractors have full border rows but different interior patterns. My expression highlighted the central cross shape and the empty corners, which uniquely matched the target. Next time, I would similarly focus on the distinctive interior pattern and symmetry to differentiate.} \\
            }
        }
    }
    & & \\ \\

    \theutterance \stepcounter{utterance}  
    & \multicolumn{4}{p{0.45\linewidth}}{
        \cellcolor[rgb]{0.95,0.95,0.95}{
            \makecell[{{p{\linewidth}}}]{
                \texttt{\tiny{[P1$|$internal]}}
                \texttt{{[}attempt 1 – valid{]}} \\
\texttt{Expression: Rows 1, 3, and 5 are fully filled with X's, while rows 2 and 4 have an X only in the first and last columns, forming a vertical cross pattern with empty cells in the middle of rows 2 and 4, and the center column (column 2) is filled in all rows except row 4, which has an X only in the first and last columns.} \\
            }
        }
    }
    & & \\ \\

    \theutterance \stepcounter{utterance}  
    & \multicolumn{4}{p{0.6\linewidth}}{
        \cellcolor[rgb]{0.9,0.9,0.9}{
            \makecell[{{p{\linewidth}}}]{
                \texttt{\tiny{[P1$\rangle$GM]}}
                \texttt{Expression: Rows 1, 3, and 5 are fully filled with X's, while rows 2 and 4 have an X only in the first and last columns, forming a vertical cross pattern with empty cells in the middle of rows 2 and 4, and the center column (column 2) is filled in all rows except row 4, which has an X only in the first and last columns.} \\
            }
        }
    }
    & & \\ \\

    \theutterance \stepcounter{utterance}  
    & & & \multicolumn{2}{p{0.3\linewidth}}{
        \cellcolor[rgb]{0.9,0.9,0.9}{
            \makecell[{{p{\linewidth}}}]{
                \texttt{\tiny{[GM$|$GM]}}
                \texttt{Expression: Rows 1, 3, and 5 are fully filled with X's, while rows 2 and 4 have an X only in the first and last columns, forming a vertical cross pattern with empty cells in the middle of rows 2 and 4, and the center column (column 2) is filled in all rows except row 4, which has an X only in the first and last columns.} \\
            }
        }
    }
    & & \\ \\

    \theutterance \stepcounter{utterance}  
    & & \multicolumn{4}{p{0.6\linewidth}}{
        \cellcolor[rgb]{0.9,0.9,0.9}{
            \makecell[{{p{\linewidth}}}]{
                \texttt{\tiny{[GM$\rangle$P2]}}
                \texttt{You are given three grids, where each of them is 5 by 5 in size.} \\
\texttt{Grids have empty cells marked with "▢" and filled cells marked with "X".} \\
\texttt{You are also given a referring expression that describes one of the given grids.} \\
\texttt{Your task is to select the grid that matches the given referring expression.} \\
\texttt{Generate only the number (in text) of the grid that matches the given expression by selecting first, second, or third.} \\
\\ 
\texttt{First:} \\
\\ 
\texttt{□ X X X X} \\
\texttt{□ X □ □ X} \\
\texttt{□ X X X X} \\
\texttt{□ □ □ □ □} \\
\texttt{□ X X X X} \\
\\ 
\texttt{Second:} \\
\\ 
\texttt{□ X X X X} \\
\texttt{□ X □ □ X} \\
\texttt{□ X X X X} \\
\texttt{□ X □ □ X} \\
\texttt{□ X X X X} \\
\\ 
\texttt{Third:} \\
\\ 
\texttt{□ X X X X} \\
\texttt{□ X □ □ X} \\
\texttt{□ X □ □ X} \\
\texttt{□ X □ □ X} \\
\texttt{□ X X X X} \\
\\ 
\texttt{Expression: Rows 1, 3, and 5 are fully filled with X's, while rows 2 and 4 have an X only in the first and last columns, forming a vertical cross pattern with empty cells in the middle of rows 2 and 4, and the center column (column 2) is filled in all rows except row 4, which has an X only in the first and last columns.} \\
\texttt{Question: Which grid does the expression refer to?} \\
\texttt{Start with the tag "Answer: ", followed by your selection. Omit any other text.} \\
            }
        }
    }
    & \\ \\

    \theutterance \stepcounter{utterance}  
    & & & \multicolumn{4}{p{0.45\linewidth}}{
        \cellcolor[rgb]{0.95,0.95,0.95}{
            \makecell[{{p{\linewidth}}}]{
                \texttt{\tiny{[P2$|$internal]}}
                \texttt{Reflections from previous episodes:} \\
\texttt{{-}} \\
\texttt{ Failed. I misinterpreted "second and fourth rows" as the internal rows (2 and 4 of the whole grid), but the expression meant rows 2 and 4 *within the hollow center*. The correct match was the first grid, which has empty internal rows 2 and 4 (of the 3x3 center). Next time, I will map the description precisely to the grid's coordinate system.} \\
\texttt{{-}} \\
\texttt{ Answer: First} \\
\\ 
\texttt{Reflection: Failed. I misinterpreted the expression's details about the cross pattern and row specifics. The description actually matched the second grid, which has a full border and a plus shape inside. I should have compared each grid more carefully to the exact pattern described, especially the "center 3x3 area has X's forming a plus shape." Next time, I will map the description directly to each grid's interior structure.} \\
\texttt{{-}} \\
\texttt{ Answer: Second} \\
\\ 
\texttt{Reflection: Succeeded. I correctly parsed the expression: border of X's, middle row full X's, second/fourth rows with X's only at first and last columns, inner three empty — matches the cross pattern in Grid 2. I compared each grid systematically, noting Grid 1 had alternating full rows, Grid 3 had a different pattern. Next time, I would continue the same method: break down the expression into specific row/column conditions and verify each grid step by step.} \\
            }
        }
    }
     \\ \\

    \theutterance \stepcounter{utterance}  
    & & & \multicolumn{4}{p{0.45\linewidth}}{
        \cellcolor[rgb]{0.95,0.95,0.95}{
            \makecell[{{p{\linewidth}}}]{
                \texttt{\tiny{[P2$|$internal]}}
                \texttt{{[}attempt 1 – valid{]}} \\
\texttt{Answer: Second} \\
            }
        }
    }
     \\ \\

    \theutterance \stepcounter{utterance}  
    & & & \multicolumn{4}{p{0.6\linewidth}}{
        \cellcolor[rgb]{0.9,0.9,0.9}{
            \makecell[{{p{\linewidth}}}]{
                \texttt{\tiny{[GM$\langle$P2]}}
                \texttt{Answer: Second} \\
            }
        }
    }
     \\ \\

    \theutterance \stepcounter{utterance}  
    & & & \multicolumn{2}{p{0.3\linewidth}}{
        \cellcolor[rgb]{0.9,0.9,0.9}{
            \makecell[{{p{\linewidth}}}]{
                \texttt{\tiny{[GM$|$GM]}}
                \texttt{second} \\
            }
        }
    }
    & & \\ \\

\end{supertabular}
}

\end{document}
